\section{Commercial heterogeneity tools: from a conditional effect to a servicing decision}
\label{exec13:commercial-heterogeneity}

A servicing manager deciding whom to contact needs more than a ranking of
predicted recoveries. An account may have a high probability of repayment
without a call, a large incremental response to a call, or a large response
whose delivery cost exceeds its value. These are different quantities.
Commercial comparison must therefore start with the population, intervention,
outcome horizon and decision rule, then ask which documented estimator matches
that question. A label such as individualized effect cannot establish that
both potential outcomes of a particular customer have become observable.
The Stata manual itself defines its individualized average treatment effect
as an average for individuals sharing the specified characteristics
\citep{exec13_ch_stata}.

The comparison here concerns a declared universe of seventeen commercial
families. Complete primary documents support a detailed comparison of Stata
and two distinct JMP pathways, together with a more limited DataRobot
weighting workflow. Saved technical material was acquired on 13 October 2026.
Its version and publication cues precede the 7 October literature cutoff,
but those cues do not authenticate every byte as it appeared on the cutoff
date. The incomplete family coverage, failed requests and unreviewed rights
are substantive limits, not evidence that a competitor lacks a feature.
No commercial product was executed for this comparison
\citep{exec13_ch_scope}.

\subsection{One population, several legitimate answers}

Let $D$ denote a binary servicing action, $Y(d)$ the recovery in dollars over
one fixed horizon under action $d$, and $G$ a risk group measured before the
action. In the generated example, lower-risk accounts have population mass
$3/4$, action probability $1/10$, effect $2$ dollars and delivery cost $1$
dollar. Higher-risk accounts have mass $1/4$, action probability $4/5$, effect
$10$ dollars and cost $12$ dollars. Potential outcomes are stipulated, with
$Y(0)=8$ and $4$ in the respective groups and $Y(1)=Y(0)+\tau_G$. Treatment
is randomized within each group. Thus conditional exchangeability and
positivity hold by construction; they are not conclusions drawn from a bank
dataset. The exact finite-population review uses thirty lower-risk and ten
higher-risk accounts, with three and eight treated respectively. This
representation realizes the declared probabilities exactly
\citep{exec13_ch_exact}.

The average effect in the eligible population is
\[
 \tau_{\mathrm{ATE}}=\tfrac34(2)+\tfrac14(10)=4.
\]
The treated fraction is $11/40$, and the average effect among treated
accounts is instead
\[
 \tau_{\mathrm{ATET}}
 =\frac{(3/4)(1/10)(2)+(1/4)(4/5)(10)}{11/40}
 =\frac{86}{11}.
\]
The larger second number reflects the concentration of treatment in the
higher-effect group. Reporting it as the gain from treating every eligible
account would change the target population. Conversely, the ATE does not
answer a question restricted to the accounts currently receiving treatment.
These differences require a target label in an interchange record, even if
two interfaces both expose a field called effect
\citep{exec13_ch_exact}.

Residualization introduces a further distinction when effect modification
is omitted from the conditioning set. Write $W$ for all nuisance controls,
$e(W)=\Pr(D=1\mid W)$ and $h(W)=\E[Y\mid W]$. If the requested effect
model is constant while the actual effect varies with $W$, the population
least-squares residual slope is
\begin{align*}
 b_{\mathrm{PO}}
 &=\frac{\E[(D-e(W))(Y-h(W))]}{\E[(D-e(W))^2]}\\
 &=\frac{\E[e(W)(1-e(W))\tau(W)]}{\E[e(W)(1-e(W))]}.
\end{align*}
The second equality follows by conditioning on $W$, using the Bernoulli
variance of $D$ and the stipulated conditional treatment law. It is an
overlap-weighted effect, which is $214/43$ in this example, rather than the
ATE of $4$. This is a target-under-misspecification calculation, not a
claim that the inspected Stata implementation silently promises an
unrestricted ATE. Its partialing-out derivation specifies an effect
$\tau(x)$ that does not additionally vary with nuisance-only controls $w$.
Changing that scientific restriction changes the interpretation of the
residual moment \citep{exec13_ch_stata,exec13_ch_exact}.

For comparison, define $\mu_d(W)=\E[Y(d)\mid W]$ and the oracle augmented
inverse-probability score
\begin{align*}
 \psi={}&\mu_1(W)-\mu_0(W)\\
 &+\frac{D\{Y-\mu_1(W)\}}{e(W)}
 -\frac{(1-D)\{Y-\mu_0(W)\}}{1-e(W)}.
\end{align*}
With the true conditional functions and exchangeability, each residual
term has conditional expectation zero. Hence $\E[\psi]=\E[\tau(W)]=4$.
Conditioning the same score on a coarser pre-action characteristic $X$
gives $\E[\psi\mid X]=\E[\tau(W)\mid X]$ when $X$ is included in the
conditioning information. Estimating those functions, fitting a conditional
mean and obtaining valid uncertainty require additional regularity and
sampling arguments. The exact score calculation alone establishes neither
finite-sample calibration nor robustness to arbitrary learner errors
\citep{exec13_ch_exact}.

This example motivates a concrete proposal requirement: record the eligible
population, treated subpopulation, effect modifiers, nuisance-only variables
and weighting target separately. A coarse risk-band request must not inherit
an estimand merely because the nuisance model uses many variables. Model
specification, target aggregation and source-population weights are part of
the scientific contract. They must survive export to a library result, a CLI
artifact and an MCP response.

\subsection{Stata: a developed conditional-effect workflow}

The inspected \texttt{cate} documentation describes continuous outcomes and
binary treatment, with partialing-out and augmented inverse-probability
weighting paths. Outcome and treatment nuisance functions can use lasso,
random forests or documented parametric alternatives; the conditional effect
can use a forest or linear regression. It distinguishes IATEs, prespecified
group average effects, and sorted groups constructed from estimated IATE
quantiles. The distinction matters operationally: a regulator-defined risk
band is a different grouping mechanism from a band discovered after fitting
the same sample \citep{exec13_ch_stata}.

The default cross-fitting specification randomly divides observations into
ten folds. Default honest forests split the tree-building sample between
split selection and leaf estimation. The manual explains that disabling
honesty disables the corresponding forest standard errors and confidence
intervals. These are useful documented controls, but their purposes should
remain distinct. Cross-fitting separates nuisance estimation from the
observation scored; honesty separates tree partitioning from leaf labeling.
Neither operation, by its name alone, demonstrates validity when accounts
repeat within households or outcomes share calendar shocks. The inspected
random-observation algorithm does not establish a household- or
time-dependent inference result; that comparison remains open rather than
being converted into a claim about every Stata command
\citep{exec13_ch_stata}.

The out-of-bag alternative requires forest outcome and treatment models.
It cannot be combined with the documented sorted-group option, and it is
an alternative to the cross-fitting setting. A proposed wrapper must
validate these combinations before launching a job. It should not expose
cross-fitting, out-of-bag prediction and arbitrary grouping as freely
composable switches. The default propensity tolerance is $10^{-5}$; the
manual describes an error for estimated propensities below that tolerance
and an option identifying offending observations. This is a documented
computational safeguard, not a certificate of scientific positivity
\citep{exec13_ch_stata}.

Missingness is similarly method-specific. The documentation permits missing
covariates when all three models use forests, whereas other configurations
omit affected observations; an explicit option forces listwise omission.
Changing learner settings can therefore change the analyzed population.
For servicing, the disappearance of accounts with incomplete contact history
may alter both estimated response and the operational population. A wrapper
needs input and retained-row counts, reasons for exclusions, and a clear
distinction between a model's missing-value routing and a causal assumption
about the missingness process \citep{exec13_ch_stata}.

Postestimation includes conditional predictions, heterogeneity tests, group
comparisons, projections and evaluation of a prespecified treatment rule.
Forest prediction standard errors use a bootstrap of little bags; the
linear-regression prediction path uses a delta method. The policy command
computes the value of a specified rule and can compare two rules. These
are concrete comparison requirements for the proposed library: a point
estimate without a target, prediction method, uncertainty convention and
policy-selection history would be a weaker deliverable. At the same time,
documentation of a prespecified-policy command is not proof that a policy
chosen to maximize the same evaluation outcomes has valid uncertainty
\citep{exec13_ch_stata}.

Stata is therefore a candidate for a licensed external reference execution,
subject to edition, automation and runtime rights. The proposal should seek
matched generated-data comparisons for ATE, conditional means and
prespecified policy value, with identical samples and effect modifiers.
It should not propose copying proprietary implementation code or equate
a CLI wrapper with validated cross-language statistical equivalence.
Exact version, seeds, defaults, retained sample and inference method would
need to accompany each result. Runtime and entitlement verification remain
open in this survey \citep{exec13_ch_scope}.

\subsection{JMP Uplift and JMP Pro Causal Treatment answer different questions}

The JMP Uplift overview describes partition models whose splits maximize
treatment differences. For a continuous outcome, candidate splits enter a
linear model through split, treatment and their interaction. For a
categorical response, the corresponding comparison uses a logistic model.
Predictor selection uses interaction significance with a Monte Carlo
adjustment for the number of candidate levels or models. This documented
adjustment addresses a predictor-selection mechanism; it is not a general
proof of causal identification or uncertainty after all possible model and
policy searches \citep{exec13_ch_jmp_uplift}.

The separately versioned JMP 16 guide makes two particularly consequential
interface details explicit. Treatment value ordering determines the
comparison, and more than two treatment levels are reduced to the first
level against all remaining levels combined. Validation data help choose
the split sequence. Those older-guide details must be rechecked in a
chosen contemporary runtime; they are not automatically a JMP 19 guarantee.
Their architectural lesson already applies: an adapter accepting three
servicing actions must not silently turn a three-action decision into a
binary first-versus-rest contrast. Likewise, validation used for selecting
a tree is not a fresh final policy-evaluation sample. A node's displayed
treatment difference or test statistic does not by itself settle that
selection issue \citep{exec13_ch_jmp_uplift}.

JMP Pro's Causal Treatment personality is a different pathway. The inspected
19.2 documentation describes binary or continuous treatment, binary or
continuous outcome, and average-effect summaries. For binary treatment,
the report distinguishes population ATE from ATET and documents
inverse-probability weighting with ratio adjustment, regression adjustment
and augmented weighting. The specified propensity model is logistic
regression, with linear or logistic outcome regression according to
outcome scale. These claims concern those documented model paths, not an
inferred native conditional-effect forest or arbitrary machine-learning
interface \citep{exec13_ch_jmp_causal}.

For continuous treatment, the prose specifies a normal conditional
treatment density with parameters fitted by linear regression. It presents
a stabilized-weight marginal structural model and a generalized propensity
score route to mean potential outcomes at specified doses. Thus a continuous
contact intensity or price change needs a dose and support definition,
not a reinterpretation of a binary-treatment coefficient. The displayed
report discusses a treatment and squared-treatment regression; its
image-rendered equations were not transcribed into this comparison.
Only inspected prose supports the implementation description
\citep{exec13_ch_jmp_causal}.

The same documentation states consistency, conditional exchangeability and
positivity as identification assumptions. Empirical overlap plots can
reveal severe support problems in observed data; they cannot establish the
absence of unmeasured confounding, a well-defined intervention or population
positivity everywhere. An attractive diagnostic plot must therefore be
stored with its interpretation and limitations. The proposal should
compare JMP's documented reports as a user-facing baseline while retaining
explicit assumption and target records \citep{exec13_ch_jmp_causal}.

The practical reuse choices differ by pathway. Uplift is a candidate
benchmark for transparent treatment-response segmentation under a suitable
experimental design. Causal Treatment is a candidate comparator for the
specified average and dose-response models. Neither choice establishes
availability of an unattended server interface, MCP service or permitted
redistribution. Product edition and rights must be reviewed before an
execution adapter is approved. A file-based result import can be designed
now, with an explicit unresolved-runtime disposition
\citep{exec13_ch_scope}.

\subsection{A documented platform workflow is a narrower claim than a native estimator}

DataRobot's official Causal AI accelerator overview describes fitting a
treatment-propensity model, calculating inverse-probability weights and
comparing outcomes in the weighted population. Its example concerns
hospital readmission; the relevant architectural point is the separation
between predictive propensity fitting and causal weighting. The overview links to a notebook whose pinned source is developed in
Section~\ref{sec:exec25-datarobot-reuse}. That study traces in-sample propensity
predictions, row/class mapping and normalized weighting, while retaining the
unresolved source-permission and native-runtime boundaries. Neither the
overview nor that source inspection establishes a native CATE engine, a
doubly robust score or calibrated inference
\citep{exec13_ch_datarobot}.

For the proposed library, the conditional adoption route is to accept
versioned nuisance predictions through a narrow data contract if their
training population, folds, target encoding and preprocessing are known.
Identification, overlap policy, target aggregation and inference remain
separate responsibilities. A prediction service becomes usable in an
orthogonal-score workflow only when its training/evaluation separation
and error conditions are justified. The current evidence does not approve
that integration or a performance comparison. It specifies what code and
rights review must precede one.

The other fourteen families are not silently dropped. BayesFusion, Bayesia,
HUGIN, Norsys, IBM, Mplus, EViews, OxMetrics, Scientific Software
International, Bayes Server and Lumina retain prior-area source leads,
but a matched heterogeneous-effect method, edition and rights record is
not established by those leads. SAS's newly saved treatment-effects PDF
is partial and excluded; failed technical routes do not settle the scope
of its products. The requested MathWorks route returned access denial,
and the causaLens request failed certificate-hostname validation. These
are acquisition limits. The survey requires successor source work before
ranking those families against the proposed capability
\citep{exec13_ch_scope}.

\subsection{Overlap, selection and the economics of a policy}

Overlap policy is not a minor numerical option. Consider a generated
population with true propensity $1/20$, $Y(1)=10$ and $Y(0)=0$. If an
unnormalized IPW expression replaces the propensity by a floor of $1/10$,
its expected treated contribution is $(1/20)10/(1/10)=5$, although the
true ATE is $10$. This deliberately narrow counterexample concerns that
unnormalized formula. It does not assert the same bias for a ratio-normalized
estimator, a correctly specified augmented estimator or a separately
defined overlap target. It shows why the adapter must preserve the exact
weighting and normalization rule, not just a boolean clipping flag
\citep{exec13_ch_exact}.

Policy value introduces costs and selection. Let $\pi(G)$ be a fixed
binary contact rule and $c(G)$ its per-contact cost. Its incremental net
value per eligible account in the constructed example is
\[
 V(\pi)=\E[\pi(G)\{\tau_G-c(G)\}].
\]
Treating no one yields zero; treating only lower-risk accounts yields
$3/4$ dollars; treating only higher-risk accounts yields $-1/2$ dollars;
treating both yields $1/4$ dollar. The highest-effect group is not the
highest-net-benefit group because its stipulated delivery cost is larger.
These are exact scenario values under fixed effects, costs and population
shares. Capacity, customer harm, uncertainty, interference and treatment
feasibility would require additional constraints before a bank decision
\citep{exec13_ch_exact}.

Even without those complications, choosing a policy on its evaluation
noise can inflate the apparent benefit. Suppose two equal true effects
are observed with independent errors taking $-1$ and $1$ with equal
probability. The expected larger reported effect is $1/2$. Selecting
with these observations and evaluating with fresh independent zero-mean
errors gives expected evaluation error zero. This elementary construction
does not accuse a vendor of using an invalid procedure. It explains why
the proposed workflow must distinguish policy training, tuning and final
evaluation, and why a prespecified-rule command cannot automatically
certify an adaptively chosen portfolio \citep{exec13_ch_exact}.

\subsection{What the proposed capability must add, reuse and leave unresolved}

The inspected alternatives make the required comparison concrete. A
credible proposed heterogeneous-effect module needs documented target
definitions and aggregation weights; nuisance and effect-model options;
reproducible splitting and missingness rules; overlap diagnostics and
failure behavior; uncertainty with its sampling-unit assumptions; and
policy-value reporting under explicit cost and selection conventions.
The commercial documentation establishes comparison requirements. It does
not demonstrate that this project's current service already implements
all of them \citep{exec13_ch_stata,exec13_ch_jmp_causal,exec13_ch_scope}.

A first implementation should use a small declared set of targets and
sampling laws. An artifact would bind the treatment contrast, outcome
horizon and currency unit; the eligible and retained populations;
pre-action effect modifiers and controls; nuisance predictions and their
fold membership; weight transformation and normalization; conditional
versus population aggregation; and the uncertainty and policy-evaluation
method. Unsupported multi-action, repeated-account or selected-policy
requests should receive a specific disposition rather than a
misinterpreted numeric result. CLI and MCP can carry this record, but
transport does not supply the absent statistical argument.

The recommended comparison sequence is to reproduce the exact target
examples, then use a prospectively specified generated design for matched
estimator and fixed-policy comparisons. Vendor execution is conditional
on a licensed, versioned runtime. Code porting requires an actual right
to reuse code; the present recommendation for proprietary methods is
documented comparison or authorized external execution. Open-source
implementation choices remain subject to their separate source, license
and numerical reviews. The immediately actionable contribution is a
scientifically explicit interface and evaluation design, with the
commercial survey still inconclusive where essential method, edition,
rights or runtime evidence is missing \citep{exec13_ch_scope}.
\label{exec13:commercial-heterogeneity-end}
